\section{The quintic upper bound}\label{sec:application}

We now apply the comparison results to the full polynomial class.
First replace $y$ by $y-F(0)$, so that $F(0)=0$. If
\[
 X(t)=-x(-t),\qquad Y(t)=y(-t),
\]
then
\[
 \dot X=Y-F(-X),\qquad \dot Y=-X.
\]
This reflection and time reversal preserve the geometric periodic
orbits and their isolation, and negate the coefficient of $x^5$.
We may therefore assume that coefficient is nonnegative. No
normalization of the other coefficients is needed.

On either open half-plane, put $u=x^2/2$ and use $y$ as coordinate
along the orbit. Equation~\eqref{eq:original-system} becomes
\begin{equation}\label{eq:quintic-profiles}
 \frac{du}{dy}=\phi_\pm(u)-y,\qquad
 \phi_\pm(u)=d_0u+b_0u^2\pm p(u),\qquad
 p(u)=e u^{1/2}+c u^{3/2}+a u^{5/2},\quad a\ge0.
\end{equation}
The plus sign corresponds to $x>0$. Indeed,
$\phi_\pm(u)=F(\pm\sqrt{2u})$; if
$F(x)=\sum_{j=1}^5 f_jx^j$, then
\[
 d_0=2f_2,\quad b_0=4f_4,\quad
 e=\sqrt2\,f_1,\quad c=2^{3/2}f_3,\quad a=2^{5/2}f_5.
\]
Thus $d_0,b_0,e,c$ are unrestricted. The profiles are smooth on
$(0,\infty)$ and continuous at zero, where they vanish.

Let $J$ be the common open interval of transverse height-zero
widths from Proposition~\ref{prop:cycle-matching}. If $J$ is empty
there is no periodic orbit. For $r\in J$, define
$\lambda_\pm(0,r),\kappa_\pm(0,r)$ by applying the inverse in
Theorem~\ref{thm:quadratic-fit} to
$(M_\pm(0,r),M_{\pm,r}(0,r),r)$.
Lemma~\ref{lem:axis-extension} makes these data interior to the
fit domain; continuity of the inverse identifies the resulting
parameters with the limits of the positive-height fits. Put
\begin{equation}\label{eq:matching-difference}
 \Delta(r)=M_+(0,r)-M_-(0,r),\qquad r\in J.
\end{equation}
Proposition~\ref{prop:cycle-matching} shows that it suffices to
bound the isolated zeros of $\Delta$ on this single interval.

\subsection{The model comparison at a matching height}

Recall the model width derivative after fitting a prescribed midpoint,
defined in Equation~\eqref{eq:fit-conditional-data}:
\[
 \widehat v(r,M,\kappa)
 =H_r\bigl(\lambda_*(r,M,\kappa),\kappa,r\bigr),
 \qquad r>0,\quad |M|<r,\quad \kappa\in\mathbb R.
\]
Implicit differentiation of the midpoint fit gives
\begin{equation}\label{eq:V-derivatives}
 \widehat v_M=\frac RP>0,\qquad
 \widehat v_\kappa=\frac GP>0.
\end{equation}
All model derivatives here are evaluated at the corresponding
parameters. For either actual profile,
\begin{equation}\label{eq:midpoint-ode}
 \frac{d}{dr}M_\pm(0,r)
 =\widehat v\bigl(r,M_\pm(0,r),\kappa_\pm(0,r)\bigr).
\end{equation}
The smooth dependence at height zero in
Lemma~\ref{lem:axis-extension}, together with the joint smooth
inverse in Theorem~\ref{thm:quadratic-fit}, justifies these identities
and the limits used below.

Lemma~\ref{lem:fitted-curvature-range} bounds each fitted curvature
between the minimum and maximum of the actual curvature along its
arc. We will use this local bound before passing to height zero.

\begin{lemma}\label{lem:isolated-zero-count}
Let $f$ be continuously differentiable on an open interval and
suppose $f'=bq$, where $b$ is continuous and positive and $q$ is
continuous and nondecreasing. Then $f$ has at most two isolated
zeros. If a continuously differentiable function has strictly
positive derivative at each of its zeros, it has at most one zero.
\end{lemma}

\begin{proof}
For the first statement, the sets where $q$ is negative, zero,
and positive occur in that order. The zero set of $q$ is an
interval, possibly empty or a singleton. Accordingly $f$ strictly
decreases, is constant, and strictly increases on the respective
sets. Each strict part has at most one zero. If the constant part
has nonzero value it contributes none. If it has positive length
and value zero, every point of it and each endpoint lying in the
domain are nonisolated zeros, and neither strict part contributes
another zero. If it is a singleton at which $f=0$, that point is
the sole minimum zero. These possibilities prove the claim.

For the second statement, suppose $r_1<r_2$ were two zeros.
Positive derivatives give $a,b$ with
$r_1<a<b<r_2$, $f(a)>0$, and $f(b)<0$. The first zero $c$ in
$[a,b]$ exists by continuity. Since $f>0$ on $[a,c)$, its
derivative at $c$, which exists, is nonpositive. This contradicts
the hypothesis.
\end{proof}

\subsection{The four coefficient cases}

\paragraph{Case 1: $e\ge0$ and $c\ge0$.}
Either $p\equiv0$, in which case the two profiles and their midpoint
functions coincide and no root in $J$ is isolated, or $p(u)>0$ for
every $u>0$. In the latter case even one periodic orbit is
impossible. To see this, suppose the two positive arches
$u_\pm(y)$ had the same two endpoints $A<B$ at height zero.
Their difference $w=u_+-u_-$ satisfies
\[
 w'=\bigl[d_0+b_0(u_++u_-)\bigr]w
       +p(u_+)+p(u_-),\qquad w(A)=w(B)=0.
\]
The coefficient is continuous on $[A,B]$, and the inhomogeneous
term is strictly positive on $(A,B)$. The integrating-factor
formula from $A$ gives $w(B)>0$, a contradiction.

\paragraph{Case 2: $e\ge0$ and $c<0$.}
Here
\begin{equation}\label{eq:third-odd}
 p'''(u)=\frac{3e-3cu+15au^2}{8u^{5/2}}>0.
\end{equation}
Proposition~\ref{prop:third-derivative} gives
$\partial_r\kappa_+(h,r)>0$ and
$\partial_r\kappa_-(h,r)<0$ for $h,r>0$. Passage to the
height-zero limits at any two specified widths in $J$ shows that
\[
 q(r)=\kappa_+(0,r)-\kappa_-(0,r)
\]
is nondecreasing on $J$.

Subtract the two instances of Equation~\eqref{eq:midpoint-ode}.
The fundamental theorem
of calculus, first in $M$ and then in $\kappa$, gives
\begin{equation}\label{eq:Delta-transport}
 \Delta'=A(r)\Delta+B(r)q,\qquad B(r)>0,
\end{equation}
with continuous coefficients. More explicitly, suppressing
$(0,r)$ in the actual parameters, one can take
\[
 \begin{split}
 A(r)&=\int_0^1
 \widehat v_M\bigl(r,M_-+\tau(M_+-M_-),\kappa_+\bigr)\,d\tau,\\
 B(r)&=\int_0^1
 \widehat v_\kappa\bigl(r,M_-,\kappa_-+\tau(\kappa_+-\kappa_-)\bigr)\,d\tau.
 \end{split}
\]
The midpoint segment remains in $(-r,r)$ and there is no
restriction on the curvature segment. Equation~\eqref{eq:V-derivatives}
therefore proves the asserted positivity.

Fix $r_*\in J$ and multiply $\Delta$ by
$\exp(-\int_{r_*}^r A(s)\,ds)$. The derivative of the resulting
function is a continuous positive factor times $q$.
Lemma~\ref{lem:isolated-zero-count} bounds its isolated zeros,
and hence those of $\Delta$, by two. This argument includes
multiple isolated roots and zero intervals.

\paragraph{Case 3: $e<0$ and $c\ge0$.}
In this case
\begin{equation}\label{eq:second-odd}
 p''(u)=\frac{-e+3cu+15au^2}{4u^{3/2}}>0,
 \qquad p''(u)\longrightarrow+\infty\quad(u\downarrow0).
\end{equation}
Fix $r_0\in J$. The peaks of the positive-height arcs at width
$r_0$ approach finite positive peaks as $h\downarrow0$, by
Lemma~\ref{lem:axis-extension}. Choose $T$ larger than both peaks
and all sufficiently nearby ones. The function $p''$ has a
positive lower bound $\delta$ on $(0,T]$: it tends to infinity
at zero and is continuous and positive away from zero.
Consequently, throughout the relevant arcs,
\[
 \phi_+''\ge2b_0+\delta,\qquad
 \phi_-''\le2b_0-\delta.
\]
Lemma~\ref{lem:fitted-curvature-range} and passage to the limit give
\begin{equation}\label{eq:curvature-gap}
 \kappa_+(0,r_0)\ge2b_0+\delta
 >2b_0-\delta\ge\kappa_-(0,r_0).
\end{equation}
At every zero of $\Delta$, Equation~\eqref{eq:midpoint-ode}
and $\widehat v_\kappa>0$ now imply $\Delta'>0$.
Lemma~\ref{lem:isolated-zero-count} gives at most one zero.

\paragraph{Case 4: $e<0$ and $c<0$.}
We have
\begin{equation}\label{eq:intercept-odd}
 up''(u)-p'(u)=
 \frac{-3e-3cu+5au^2}{4\sqrt u}>0.
\end{equation}
Writing $k_\pm=\phi_\pm'$, Equation~\eqref{eq:intercept-odd}
is precisely
\[
 \left(\frac{k_+(h)-d_0}{h}\right)'>0,\qquad
 \left(\frac{k_-(h)-d_0}{h}\right)'<0.
\]
Proposition~\ref{prop:intercept} therefore gives
\begin{equation}\label{eq:positive-height-intercepts}
 \lambda_+(h,r)-h\kappa_+(h,r)<d_0,\qquad
 \lambda_-(h,r)-h\kappa_-(h,r)>d_0
 \quad(h,r>0).
\end{equation}
The strict inequalities persist at height zero, but this requires
more than simply taking limits. Fix a transverse peak of the plus
profile, and let $r(h)$ be its width at height $h$, with
$r(h)\to r_0\in J$. The endpoint data converge to an interior
triple $(M_0,v_0,r_0)$, with $|M_0|<r_0$ and $|v_0|<1$.
The joint smooth inverse in Theorem~\ref{thm:quadratic-fit}
shows that $\lambda(h,r(h))$ and $\kappa(h,r(h))$ have finite
limits. Put
\[
 I(h)=\lambda(h,r(h))-h\kappa(h,r(h))-d_0.
\]
The derivative along this fixed-peak characteristic is
\begin{equation}\label{eq:strict-boundary-intercept}
 \frac{dI}{dh}=DI=\frac{S+hR}{G}\bigl(\lambda-k_+(h)\bigr).
\end{equation}
Since $e<0$, $k_+(h)\to-\infty$, whereas $\lambda$ remains
bounded. Thus $dI/dh>0$ for all sufficiently small positive $h$.
Choose one such $h_1$. Equation~\eqref{eq:positive-height-intercepts}
and integration in decreasing $h$ give
\[
 I(h)\le I(h_1)<0\qquad(0<h<h_1).
\]
Taking the limit yields $\lambda_+(0,r_0)<d_0$. On the minus side,
$k_-(h)\to+\infty$, so the same argument with all signs reversed
gives $\lambda_-(0,r_0)>d_0$. Therefore, throughout $J$,
\begin{equation}\label{eq:strict-slope-gap}
 \lambda_+(0,r)<d_0<\lambda_-(0,r).
\end{equation}

At a zero of $\Delta$, the two model midpoints agree. Because
$H_\lambda,H_\kappa>0$, Equation~\eqref{eq:strict-slope-gap}
forces $\kappa_+(0,r)>\kappa_-(0,r)$: otherwise increasing
$\lambda_+$ to $\lambda_-$ and then $\kappa_+$ to $\kappa_-$
would strictly increase $H$. Equation~\eqref{eq:midpoint-ode}
again gives $\Delta'>0$ at every zero, and there is at most one.

These four cases exhaust all $e,c\in\R$ with $a\ge0$, including
zero leading coefficients. By Proposition~\ref{prop:cycle-matching},
they prove the upper bound in Theorem~\ref{thm:main} for every
allowed polynomial, without a restriction on amplitude or on
the multiplicity of a limit cycle.
